\section{排序与分治}

本章收录以\textbf{排序与分治}为核心思想的经典问题模板。这类问题不依附于某一数据结构：分治解法依托递归结构在合并时统计贡献，数据结构解法（如树状数组）则依托值域统计，两者常可互相印证。

\subsection{逆序对 (Inversion Count)}

\textbf{用途：}
求序列中满足 $i < j$ 且 $a_i > a_j$ 的二元组 $(i, j)$ 个数，即逆序对数量。经典计数问题，有两种 $O(n \log n)$ 标准解法，核心都是「维护已扫描部分、逐个统计贡献」，按数据特点选用：

\begin{itemize}
  \item \textbf{树状数组法：} 以值域为下标做单点加与前缀和，需要\textbf{离散化}；统计思想通用，可直接推广到二维偏序、动态插入等场景。
  \item \textbf{归并排序法：} 借助分治结构在合并时一次统计跨段贡献，\textbf{无需离散化}、常数更小，但仅适用于静态序列的一次性计数。
  \item \textbf{溢出提醒：} 逆序对最多 $\dfrac{n(n-1)}{2}$，两种方法都必须用 \code{long long} 计数。
\end{itemize}

\subsubsection{归并排序求逆序对 (By Merge Sort)}

\textbf{思路：}
借助归并排序的分治结构，在合并两个有序段的同时统计跨越左右的逆序对，无需离散化。

\textbf{核心概念：}
\begin{itemize}
  \item \textbf{分治计数：} 逆序对按位置分三类：都在左半段、都在右半段、跨越左右。前两类在递归求解子问题时自动统计，第三类在合并时统计。
  \item \textbf{合并时计数：} 左右两段各自有序，当需要取出右段元素 $a_j$ 时，左段剩余元素 $a_i \sim a_{\text{mid}}$ 全部大于 $a_j$，一次性贡献 $\text{mid} - i + 1$ 对。
  \item \textbf{稳定性：} 相等时优先取左段（比较用 $a_i \le a_j$），保证相等元素不计入逆序对，排序结果也保持稳定。
\end{itemize}

\textbf{时间复杂度：} $O(n \log n)$，空间复杂度 $O(n)$。

\lstinputlisting[language=C++]{codes/03-sorting-01.cpp}

\subsubsection{树状数组求逆序对 (By Fenwick Tree)}

\textbf{思路：}
从右往左扫描序列，树状数组维护「已出现过的值」的个数，对每个元素累加右侧比它小的元素个数。

\textbf{核心概念：}
\begin{itemize}
  \item \textbf{离散化：} 值域可能很大（如 $10^9$）或含负数，先把所有值排序去重，映射到 $1 \sim n$ 的排名，再以排名为下标使用树状数组。
  \item \textbf{单点加与前缀和：} \code{modify(rk[i])} 将当前值的计数加一；\code{query(rk[i] - 1)} 求严格小于 $a_i$ 的元素个数，恰好排除相等的元素（相等不计入逆序对）。
  \item \textbf{对偶思路：} 也可从左往右扫描，累加「已插入元素中大于 $a_i$ 的个数」，即 $\text{query}(n) - \text{query}(rk_i)$。
\end{itemize}

\textbf{时间复杂度：} 离散化与扫描均为 $O(n \log n)$，空间复杂度 $O(n)$。

\lstinputlisting[language=C++]{codes/03-sorting-02.cpp}
